\documentclass[a4paper,11pt]{article}
\usepackage[T2A]{fontenc}
\usepackage[utf8]{inputenc}
\usepackage[russian]{babel}
\usepackage{amsmath,amssymb,mathtools}
\usepackage{cmbright}
\renewcommand{\familydefault}{\sfdefault}
\usepackage{icomma}
\usepackage[a4paper,left=16mm,right=16mm,top=12mm,bottom=11mm]{geometry}
\usepackage{microtype}
\usepackage{multicol}
\usepackage{xcolor}
\usepackage{hyperref}
\definecolor{accent}{HTML}{15747A}
\definecolor{darkaccent}{HTML}{164C55}
\definecolor{textgray}{HTML}{283238}
\definecolor{muted}{HTML}{65747A}
\definecolor{rulegray}{HTML}{CADADC}
\color{textgray}
\hypersetup{hidelinks}
\setlength{\parindent}{0pt}
\setlength{\parskip}{0pt}
\setlength{\columnsep}{8mm}
\setlength{\abovedisplayskip}{1mm}
\setlength{\belowdisplayskip}{1mm}
\emergencystretch=2em
\newcommand{\sect}[1]{\par\vspace{2.2mm}{\normalsize\bfseries\color{darkaccent}#1}\par\vspace{0.3mm}{\color{rulegray}\rule{\linewidth}{0.45pt}}\par\vspace{0.8mm}}
\newcommand{\fml}[1]{\par\noindent{\normalsize\(\displaystyle #1\)}\par\vspace{2.8mm}}
\begin{document}
\thispagestyle{empty}

\hypersetup{pdftitle={Главные формулы. Григорий Архипов},pdfauthor={Лёвин Артём Александрович}}
\begin{center}
{\small\bfseries\color{accent}УЧЕБНАЯ ПАМЯТКА \quad\textbullet\quad ОСНОВНЫЕ ФОРМУЛЫ}\par
\vspace{1mm}
{\fontsize{21}{25}\selectfont\bfseries\color{darkaccent} Многочлены и прикладные формулы\par}
\vspace{0.8mm}
{\small\color{muted}Григорий Архипов\quad\textbullet\quad ЕГЭ · профильная математика\quad\textbullet\quad 12.10.2026}
\end{center}
\vspace{2mm}
\noindent{\color{accent}\rule{\linewidth}{0.9pt}}
\vspace{1.3mm}
{\bfseries\color{darkaccent}Главные принципы решения}\par
\vspace{1mm}
\begin{multicols}{2}
\raggedcolumns
\small
\noindent\textbullet\quad Перед делением многочлена проверяйте потенциальный корень подстановкой.\par\vspace{0.75mm}
\noindent\textbullet\quad При схеме Горнера не пропускайте нулевые коэффициенты отсутствующих степеней.\par\vspace{0.75mm}
\columnbreak
\noindent\textbullet\quad В прикладной задаче указывайте допустимый диапазон величин.\par\vspace{0.75mm}
\noindent\textbullet\quad Корни квадратного уравнения сверяйте с ограничениями задачи.\par\vspace{0.75mm}
\end{multicols}
\vspace{1.1mm}
\noindent{\color{accent}\rule{\linewidth}{0.5pt}}
\vspace{0.6mm}
\begin{multicols}{2}
\raggedcolumns
\normalsize
\sect{Теорема Безу}
\fml{P(a)=0\Leftrightarrow P(x)=(x-a)Q(x)}
\fml{P(x)=(x-a)Q(x)+P(a)}
\sect{Квадратное уравнение}
\fml{ax^2+bx+c=0\quad(a\ne0)}
\fml{D=b^2-4ac}
\fml{x_{1,2}=\frac{-b\pm\sqrt D}{2a}}
\columnbreak
\sect{Формулы сокращённого умножения}
\fml{(a+b)(a-b)=a^2-b^2}
\fml{(a\pm b)^2=a^2\pm2ab+b^2}
\sect{Проценты и модель}
\fml{k=1\pm\frac p{100}}
\fml{S_{\text{нов}}=kS_{\text{исх}}}
\end{multicols}
\vfill
{\color{accent}\rule{\linewidth}{0.8pt}}\par\vspace{1.5mm}
{\small\color{muted}Лёвин Артём Александрович эксклюзивно для Григория Архипова}
\end{document}